\chapter{Introduction}
\label{chap:introduction}

\section{Motivation: Decoding Affect from Neurophysiology as an Electrical Engineering Problem}

This dissertation is positioned first and foremost as an electrical engineering study of structured, noisy, nonstationary biosignals. The application domain is affective neurophysiology, but the central technical problem is one of signal representation, nonlinear state transformation, stochastic dynamical processing, and graph-structured analysis under severe uncertainty. In this setting, scalp electroencephalography (EEG) is treated not as a direct readout of latent mental state, but as a multichannel, volume-conducted, low-amplitude observation of distributed cortical activity, contaminated by measurement noise, physiological artifacts, inter-subject variability, and nonstationary background dynamics. The dissertation therefore addresses a classical engineering challenge: how to construct a signal processing framework that transforms these observations into a representation in which affective and diagnostically informative structure is more observable, more separable, and---critically---more \emph{interpretable}. A clinician cannot act on a classification output from an opaque model. The architecture must produce not only predictions but traceable, clinically meaningful explanations of how and why neural response organization differs across individuals, affective conditions, and diagnostic categories.

The governing question that organizes the investigation is this: \emph{can hidden dynamical order be recovered from noisy neural measurements through nonlinear transformation, and under what conditions does that order remain stable, compressible, structurally coherent, and interpretable in terms that are meaningful to clinicians?} Raw scalp EEG is a corrupted surface observation of distributed cortical dynamics. The central conjecture of this dissertation is that a staged neuromorphic architecture---comprising a spiking reservoir, a temporal coding bottleneck, and a graph-organized spatial analysis---can convert those corrupted observations into a representation where latent temporal, geometrical, and relational order becomes observable, separable, and traceable back to the input signal through an explicit, inspectable signal chain. Classification accuracy is one instrument for evaluating that recovery, but it is not the sole or even the primary measure of success. The deeper measures are whether the staged representation exposes forms of signal organization---operating regimes, compression boundaries, relational regime limits, dynamical response properties, and cross-scale correspondences---that no single-stage approach can access, and whether each output variable can be traced to specific properties of the input ERP waveform without reliance on post-hoc explanation methods.

\subsection{The Engineering Difficulty of Scalp EEG}

Scalp EEG records voltage fluctuations on the order of $1$--$100\;\mu$V, arising from the superposition of post-synaptic potentials in cortical pyramidal cells. These signals are attenuated and spatially blurred by volume conduction through cerebrospinal fluid, skull, and scalp before reaching the electrode, producing a severely underdetermined inverse problem: the mapping from cortical source activity to scalp observation is many-to-one and ill-conditioned \cite{nunez_srinivasan_2006}. The effective spatial resolution of high-density scalp EEG is on the order of several centimeters, orders of magnitude coarser than functional magnetic resonance imaging (fMRI), though EEG compensates with millisecond temporal resolution that fMRI cannot approach.

The signal is further corrupted by artifacts that are often larger in amplitude than the neural signals of interest. Ocular artifacts from blinks and saccades can exceed $100\;\mu$V, dwarfing the $1$--$10\;\mu$V event-related potentials (ERPs) that index cognitive and affective processing. Muscular contamination, cardiac interference, line noise, and electrode impedance drift all contribute to a measurement environment in which the quantity of interest is deeply embedded in nuisance variability. Across trials within a single subject, EEG responses to identical stimuli exhibit substantial amplitude and latency jitter. Across subjects, anatomical differences in cortical folding, skull thickness, and electrode placement create systematic inter-individual variability that compounds the within-subject noise.

Most critically, EEG is \emph{nonstationary}: the statistical properties of the signal evolve continuously over time. Background oscillatory activity shifts in power and frequency as a function of arousal, attention, fatigue, and task engagement. Event-related responses are transient and temporally localized, lasting tens to hundreds of milliseconds. Any inference architecture that treats EEG epochs as draws from a fixed distribution ignores this fundamental characteristic of the data-generating process.

\subsection{Clinical Motivation: The Biomarker Translation Gap}

The urgency of developing new computational methods for EEG is motivated by a fundamental gap in clinical electrophysiological biomarker research. Conditions such as Major Depressive Disorder (MDD), Post-Traumatic Stress Disorder (PTSD), Substance Use Disorder (SUD), and Generalized Anxiety Disorder (GAD) are associated with measurable electrophysiological alterations that have been reliably demonstrated at the group level \cite{gehring_erp_2012, nelson_hajcak_2017}. Translating these group-level findings into clinically useful individual-level predictors remains an open scientific problem. A primary reason is that standard protocols emphasize within-person contrasts and trial-averaged waveforms, which suppress the between-person variance that diagnostic and prognostic applications require \cite{weinberg_hajcak_2011}.

A second reason is that psychiatric diagnosis is transdiagnostic in nature: the same electrophysiological alteration may manifest across multiple disorders, and the same disorder may produce different electrophysiological signatures depending on comorbidity, medication status, and symptom profile. A computational framework that targets a single diagnostic boundary (e.g., MDD vs.\ healthy control) captures only one dimension of the distributed, comorbidity-modulated, and condition-dependent structure that characterizes clinical EEG data. The framework developed in this dissertation embraces this complexity: the goal is to build a signal processing pipeline that produces multi-scale representations amenable to both classification and clinical network characterization across the transdiagnostic spectrum.

\subsection{The SHAPE Study as the Principal Validation Setting}

The primary real-data validation target is the SHAPE (Stony Brook Health and Psychopathology Evaluation) Community study, a clinical EEG investigation designed to examine affective processing across the full spectrum of psychopathology. SHAPE collects multichannel scalp EEG from $211$ participants---including individuals with MDD, PTSD, SUD, GAD, ADHD, eating disorders, and healthy controls---while they view emotionally evocative images from the International Affective Picture System (IAPS) \cite{lang_iaps_2008} in three conditions: negative, neutral, and pleasant. The study provides per-subject, per-condition, trial-averaged ERP matrices together with comprehensive clinical metadata including over $40$ binary diagnostic variables, medication status, and demographics.

SHAPE is well suited to the present framework for several reasons. First, its trial-averaged ERP structure is compatible with the temporal integration window of the spiking reservoir. Because SHAPE provides trial-averaged ERP matrices as the principal validation input, the present dissertation validates spike-based temporal transformation on evoked affective structure rather than on full single-trial asynchronous variability; extending the framework to single-trial EEG is an explicit direction for future work. Second, its transdiagnostic sample (mean comorbidity $2.8$ diagnoses per subject) enables clinical network analysis that goes beyond any single diagnostic boundary. Third, its multichannel EEG ($34$ channels) provides the spatial structure required for graph-based analysis. Fourth, the three-condition design provides a non-trivial classification task that directly probes the pipeline's ability to extract emotion-related temporal and spatial structure.

\subsection{Why EEG? An Engineering Justification}

The motivation for using EEG is engineering as much as physiological. EEG provides millisecond temporal resolution, natural multichannel structure, and a direct connection to transient cortical dynamics. At the same time, it presents exactly the difficulties that motivate new methodology: low signal-to-noise ratio, nonstationarity, cross-subject heterogeneity, uncertain spatial mixing, and partially observable latent dynamics. Any architecture that succeeds under these conditions must simultaneously address temporal coding, robustness to noise, dimensionality reduction, and relational structure across channels.

The engineering scope of scalp EEG is well defined. It provides millisecond temporal resolution at the cost of spatial precision relative to fMRI, source-level specificity relative to intracranial recordings, and metabolic information relative to positron emission tomography. The dissertation's signal-processing contributions operate within the information content of the scalp-recorded signal, and downstream clinical interpretation respects these measurement-model boundaries.

\subsection{Formal Problem Statement: Affective EEG as an Inverse Problem}
\label{sec:inverse_problem}

The preceding discussion motivates a precise mathematical formulation of the inference problem that ARSPI-Net is designed to address. This formulation unifies the engineering, clinical, and theoretical threads of the dissertation into a single inverse-problem framework.

Let $\mathcal{P}$ denote the population of individuals, each indexed by $p \in \mathcal{P}$. For every individual~$p$:
\begin{itemize}
    \item Let $D_p$ denote the unknown internal brain dynamics---a high-dimensional, nonlinear, non-stationary dynamical system whose state evolves according to some operator $\mathcal{F}$.
    \item Let $\delta_p \in \Delta$ denote the clinical diagnostic profile (e.g., $\Delta = \{\text{MDD}, \text{SUD}, \text{PTSD}, \text{GAD}, \text{ADHD}, \ldots\}$), which modulates $D_p$.
\end{itemize}

A known affective stimulus $s \in S$ (an IAPS image characterized by arousal and valence parameters) is presented. The internal neural response is generated by the stimulus-driven evolution
\begin{equation}
    r_p = \mathcal{F}(D_p, s; \delta_p),
    \label{eq:forward_dynamics}
\end{equation}
where $\mathcal{F}$ encodes the dynamics of cortical activity under stimulus drive, modulated by the individual's diagnostic profile.

The experimenter observes only the scalp EEG/ERP signal
\begin{equation}
    x_p = g(r_p) + \eta,
    \label{eq:observation}
\end{equation}
where $g$ is the forward volume-conduction operator (mapping cortical sources to scalp potentials) and $\eta$ encompasses measurement noise and physiological artifacts. In the SHAPE dataset, trial averaging and standard preprocessing have already been applied, so the observed $x_p(s) \in \mathbb{R}^{T \times C}$ ($T$ time samples, $C = 34$ channels) represents the scalp-level ERP matrix for individual~$p$ under stimulus condition~$s$. The volume-conduction operator $g$ is embedded in the observation; the dissertation works with the scalp-level signal directly rather than attempting cortical source reconstruction.

The central scientific question of this dissertation, stated formally, is:
\begin{equation}
    \text{Recover structured properties of } D_p \text{ from the set } \{(s,\, x_p(s))\}_{s \in S}.
    \label{eq:inverse_problem}
\end{equation}
Equivalently, the goal is to construct a mapping
\begin{equation}
    \phi: x_p(s) \;\mapsto\; \hat{D}_p
    \label{eq:phi}
\end{equation}
such that $\phi$ satisfies three requirements:
\begin{enumerate}
    \item \textbf{Stability}: $\phi$ is robust under measurement noise, physiological artifacts, and inter-subject variability in the forward operator $g$.
    \item \textbf{Interpretability}: the internal variables of $\phi$ map onto measurable ERP components (latency, amplitude, topography) and dynamical quantities (firing rate, temporal complexity, autocorrelation structure) that are meaningful to both engineers and clinicians.
    \item \textbf{Clinical generalizability}: $\phi$ generalizes across the heterogeneous transdiagnostic population of the SHAPE cohort, and the recovered properties of $\hat{D}_p$ are informative about diagnostic modulation---i.e., the relationship between $\hat{D}_p$ and $\delta_p$ is empirically testable.
\end{enumerate}

ARSPI-Net instantiates $\phi$ as a staged pipeline: the spiking reservoir performs a nonlinear temporal expansion of $x_p(s)$, the BSC$_6$-PCA-64 coding bottleneck compresses the resulting spike representation into a structured embedding, and the graph-topological and dynamical analysis layers extract spatial and temporal properties of $\hat{D}_p$ across three representational scales. The affective condition labels ($s \in \{\text{Negative, Neutral, Pleasant}\}$) serve as controlled probes for evaluating the quality of the recovered structure; the clinical diagnostic metadata ($\delta_p$) serve as the downstream validation target for whether the recovered organization is diagnostically informative.

Each subsequent chapter addresses a specific aspect of the $\phi$ construction and its three requirements. Chapter~\ref{chap:lsm} establishes the reservoir's operating regime (stability). Chapter~\ref{chap:embeddings} characterizes the coding bottleneck (stability and compression). Chapter~\ref{chap:arspi_net} evaluates the embedding geometry and the graph-propagation regime on real clinical data (interpretability and generalizability). Chapter~\ref{chap:dynamics} extracts dynamical descriptors of $\hat{D}_p$ (interpretability). Chapter~\ref{chap:synthesis} tests whether temporal and spatial properties of $\hat{D}_p$ are systematically coupled (structural coherence). Chapter~\ref{chap:conclusion} consolidates the boundary conditions under which $\phi$ succeeds and identifies the regimes where it does not.


\section{The Representation Gap and the Neuromorphic Opportunity}

\subsection{Motivating an Alternative to Conventional Deep Learning for Biosignals}

Architectures such as LSTMs and Transformers have demonstrated substantial success in generic time-series modeling \cite{ismail_fawaz_deep_2019}. When applied to physiological signals, however, several structural properties of these architectures motivate the investigation of alternative paradigms.

First, their dominant computational primitive is dense matrix multiplication, which is expensive in energy and memory traffic. The computational cost of deep learning is growing unsustainably: data centers consumed an estimated $1$--$1.5\%$ of global electricity in 2023, with AI workloads growing exponentially. In healthcare, where inference must eventually occur at the edge---on wearable devices, in ambulances, in resource-limited clinics---the assumption that GPU clusters are always available is untenable.

Second, the nonstationary character of EEG complicates the stationarity assumptions under which conventional models are easiest to train. Third, these models do not explicitly exploit the event-like, transient, and dynamically sparse structure present in neuroelectric data. Fourth, their hidden representations are difficult to interpret in systems-theoretic terms, which is particularly important when the target application is clinically meaningful inference.

These observations motivate the investigation of a fundamentally different computational paradigm. The present dissertation provides empirical evidence for the value of one such alternative: neuromorphic spiking computation.

\subsection{The Interpretability Imperative in Clinical Neuroscience}

The fourth observation above---that conventional deep learning representations are difficult to interpret in systems-theoretic terms---deserves special emphasis, because it is the observation most directly relevant to clinical translation.

In clinical neuroscience, a model's prediction is only as valuable as its explanation. A psychiatrist evaluating whether a patient's ERP profile is consistent with Major Depressive Disorder cannot act on ``the model predicts MDD with $87\%$ confidence.'' The clinician needs to know \emph{which} temporal features of the neural response are atypical, \emph{how} the patient's electrode-level dynamics differ from the normative population, and \emph{whether} the spatial organization of the patient's brain network shows patterns associated with the disorder. These are interpretability requirements that no post-hoc explanation method---not SHAP values, not attention rollout, not gradient-based saliency---can fully satisfy when applied to a model whose internal representations are learned end-to-end and bear no guaranteed relationship to the input signal's temporal or spatial structure.

The distinction is between \emph{post-hoc interpretability} and \emph{intrinsic interpretability}. Post-hoc methods attempt to explain a model's decision after the fact, by perturbing inputs or attributing importance to features. These methods are valuable but fundamentally limited: they explain what the model attended to, not what the signal contains. Intrinsic interpretability, by contrast, means that the model's internal variables are themselves meaningful quantities---that the intermediate representations map directly onto named, measurable properties of the input signal. A model with intrinsic interpretability does not need to be explained; it \emph{is} the explanation.

ARSPI-Net is designed for intrinsic interpretability. Its staged pipeline produces a sequence of named, traceable quantities at each processing stage: membrane potential trajectories and spike rasters at the reservoir stage; binned spike counts aligned to ERP temporal windows at the coding stage; per-channel embedding vectors at the compression stage; seven named dynamical descriptors (firing rate, rate entropy, permutation entropy, autocorrelation time, active ratio, spike-count variance, mean inter-spike interval) at the characterization stage; and graph-topological metrics (weighted strength, clustering coefficient, betweenness centrality) at the spatial stage. Every one of these quantities has a physical unit, a direct mapping to the input ERP waveform, and an interpretation that a clinician can evaluate. This is the engineering contribution that no black-box architecture provides, and it is the primary reason ARSPI-Net exists.

\subsection{Spiking Neural Networks: A Different Computational Paradigm}

Spiking Neural Networks (SNNs) offer a fundamentally different computational paradigm \cite{maass_networks_1997}. Rather than propagating continuous activations through densely synchronized layers, SNNs operate through sparse, asynchronous spikes generated by internal dynamical states. This creates two engineering opportunities.

The first is \emph{temporal expressivity}: information can be encoded in spike timing, synchrony, bursting structure, and transient state evolution, rather than only in static amplitude vectors. For a signal like EEG, whose information content is concentrated in the timing and morphology of brief transient events, this representational format is a principled and engineering-relevant match.

The second is \emph{potential deployment efficiency}: on dedicated neuromorphic hardware, each synaptic operation consumes approximately $20$~pJ compared to approximately $900$~pJ for a 32-bit multiply-accumulate on conventional CMOS---a $45\times$ advantage per operation \cite{davies2018loihi}. The trend toward neuromorphic hardware for energy-constrained deployment is well established and accelerating \cite{schuman2022neuromorphic}. This dissertation positions its framework to be compatible with that trajectory, although the energy advantage is motivational and is not experimentally measured in the present software implementation.

Among SNN models, the Liquid State Machine (LSM) is particularly attractive because it separates nonlinear temporal expansion from downstream readout. A fixed recurrent reservoir transforms the input into a high-dimensional state trajectory, and the downstream learner operates on that transformed state. Because the reservoir weights are not trained, the internal dynamics remain analyzable as a driven nonlinear system---a property exploited in Chapter~\ref{chap:dynamics} to extract interpretable dynamical biomarkers.

\subsection{The Graph Question: When Does Inter-Channel Structure Help?}

The LSM alone does not address the second key feature of EEG: multichannel relational structure across electrodes. Each electrode captures a different spatial mixture of underlying cortical activity, and the relationships between channels carry information about functional organization. This motivates coupling the reservoir to a graph-based analysis layer.

However, the theoretical properties of graph-based processing on EEG electrode arrays have not been systematically investigated. Graph neural networks based on message-passing---the iterative averaging of neighbor node features---have become the dominant paradigm for multichannel EEG analysis \cite{Zhong2020EEGEmotionGCN, Song2018EEGEmotionDGCNN}. The field has adopted this paradigm without characterizing its behavior on the small, dense electrode graphs typical of EEG ($\leq 64$ channels, $> 10\%$ edge density, $< 1000$ training samples). On such graphs, the smoothing inherent in message-passing may reduce inter-channel contrasts that carry discriminative information---a question the EEG-GNN literature has not yet explored.

A central contribution of this dissertation is to investigate this question empirically and to characterize the graph propagation regime for clinical EEG. The graph component of ARSPI-Net is therefore investigated as \emph{both} a candidate classification mechanism and an interpretability tool. The dissertation treats the question of which graph operator is appropriate for this data regime as an open empirical and theoretical problem whose answer constitutes a novel scientific contribution.


\section{Research Objectives and Contributions}

This dissertation introduces \emph{Affective Reservoir-Spike Processing and Inference Network (ARSPI-Net)}, a hybrid neuromorphic framework for extracting structured, interpretable, multi-scale representations from clinical EEG \cite{lane2023towards, lane2024towards}. ARSPI-Net instantiates the mapping $\phi$ defined in Section~\ref{sec:inverse_problem}: it transforms noisy scalp-level ERP observations into a staged representation from which structured properties of the unknown internal dynamics $D_p$ can be recovered, characterized, and related to clinical diagnostic status $\delta_p$---with every intermediate variable traceable to the input signal through an explicit, inspectable signal chain. The central thesis is that this staged neuromorphic architecture---comprising a spiking reservoir with no learned recurrent weights, a temporal coding bottleneck, and a graph-organized spatial analysis---produces a representation framework from which both discriminative classification and intrinsically interpretable clinical characterization can be extracted, and that a rigorous characterization of each component's operating regime constitutes a scientific contribution to the field. The architecture's primary value is not peak classification accuracy---which belongs to conventional trained models---but the combination of geometric transparency, intrinsic interpretability, and clinically meaningful structural characterization that no black-box model can provide. To our knowledge, no existing EEG analysis system achieves intrinsic interpretability at all four levels that ARSPI-Net targets: temporal traceability (the ability to trace outputs to specific ERP temporal windows), geometric transparency (quantitative decomposition of embedding structure), dynamical characterization (named descriptors with physical units that predict classical ERP scalars), and systems-level correspondence (cross-scale coupling between temporal and spatial layers). ProtoPMed-EEG \cite{barnett_protopmed_2024}---the current gold standard for clinically validated interpretable EEG---achieves prototype-based transparency but not dynamical characterization or coupling analysis. NeuCube \cite{kasabov_neucube_2014}---the most mature SNN architecture for brain data---achieves temporal alignment and connectivity analysis but does not formally decompose embedding geometry or measure cross-scale coupling. EEGNet with SHAP provides post-hoc approximations but satisfies no level intrinsically. ARSPI-Net is designed to be, to our knowledge, the first architecture that achieves all four levels through a single staged pipeline, with experimental validation at each level across Chapters~\ref{chap:embeddings}--\ref{chap:synthesis}.

\subsection{Research Objectives}

The specific research objectives are:
\begin{enumerate}
    \item To design and validate a Liquid State Machine as a nonlinear temporal feature extractor for EEG, first under controlled component-level experiments and then on real affective EEG data from the SHAPE study ($N = 211$ subjects).
    \item To develop a principled pipeline for transforming high-dimensional spike activity into compact, robust, and discriminative embeddings, including coding-scheme comparison, feature-space analysis, and dimensionality reduction.
    \item To systematically evaluate how inter-channel graph structure should be exploited---testing message-passing graph neural networks, relational feature engineering, and graph-topological analysis---and to characterize the architectural and data-regime conditions that govern each approach's operating properties.
    \item To characterize the internal driven reservoir dynamics using interpretable dynamical biomarkers so that classification performance is complemented by an engineering description of how the reservoir transforms its inputs.
    \item To synthesize the temporal-dynamical and graph-topological views of the system into a unified structure--function framework, testing whether local reservoir dynamics and global network organization are systematically related.
\end{enumerate}

\subsection{Specific Contributions}

The dissertation makes the following specific contributions:
\begin{enumerate}
    \item \textbf{An intrinsically interpretable neuromorphic measurement framework for clinical EEG.} ARSPI-Net produces a staged pipeline in which every intermediate variable---membrane trajectories, spike rasters, binned temporal codes, per-channel embeddings, dynamical descriptors, graph-topological metrics, and structure--function coupling indices---is a named, physically grounded quantity that traces directly to the input ERP waveform. This intrinsic interpretability distinguishes ARSPI-Net from conventional deep learning architectures (EEGNet, LSTM, GRU) whose internal representations are learned end-to-end and require post-hoc explanation methods. The architecture's primary value is not peak classification accuracy but clinically actionable structural characterization.
    \item \textbf{A geometric characterization of subject-covariance entanglement in affective EEG.} Variance decomposition of the ARSPI-Net embedding reveals that subject identity accounts for $62.6\%$ of embedding variance versus $8.7\%$ for emotional condition ($\rho = 7.2$). Per-subject centering recovers $+13$ to $+21$ percentage points of buried accuracy across all seven tested representations, including EEGNet ($72.0\% \to 89.1\%$), raw EEG ($70.5\% \to 88.4\%$), GRU ($59.9\% \to 78.4\%$), and ARSPI-Net ($59.4\% \to 78.8\%$). PCA-based nuisance regression is provably lossy because subject and condition share a principal subspace. This discovery---enabled by ARSPI-Net's geometric transparency---applies to every model in the comparison and motivates subject-covariance correction as a preprocessing principle for all affective EEG systems.
    \item \textbf{A graph-propagation regime characterization for small dense electrode graphs.} The Propagation Operating Characteristic (Chapter~\ref{chap:arspi_net}) demonstrates that no graph operator---smoothing, high-pass, residual, difference-aware, or attention-based---exceeds the non-propagated baseline on the $34$-channel, $20\%$-density electrode graph. This identifies a novel regime boundary with a quantitative design criterion for the EEG-GNN field. Separately, graph-topological analysis yields transdiagnostic network phenotypes, with the strongest signal in Substance Use Disorder ($p = 0.0004$).
    \item \textbf{An interpretable dynamical characterization framework.} Chapter~\ref{chap:dynamics} develops seven named dynamical descriptors---each with a physical unit, a direct mapping to ERP temporal structure, and a clinically meaningful interpretation---and validates them through echo-state verification, cross-initialization reliability, and surrogate testing. The descriptors organize subjects along a continuous excitability--persistence axis that is condition-modulated (Friedman $p = 6 \times 10^{-6}$) and not reducible to psychiatric diagnosis. This provides clinicians with a one-dimensional characterization of each patient's reservoir response organization.
    \item \textbf{A structure--function coupling analysis.} Chapter~\ref{chap:synthesis} unifies the temporal-dynamical and graph-topological characterizations, testing whether local reservoir dynamics and global network topology are systematically related---a direct test of cross-scale structural coherence in the ARSPI-Net representation.
\end{enumerate}

\subsection{Dissertation Organization}

The dissertation is organized to make the logic of ARSPI-Net's development cumulative.

Chapter~\ref{chap:background} reviews the signal-processing, dynamical-systems, neuromorphic, and graph-learning background, including the theoretical foundations of reservoir computing, graph signal processing, and the limits of message-passing on small dense graphs.

Chapter~\ref{chap:lsm} establishes the reservoir as a controlled temporal processing module through four systematic experimental modules, identifying the parameter regime carried forward into the embedding and real-EEG chapters.

Chapter~\ref{chap:embeddings} develops the spike-based embedding pipeline, specifies the BSC$_6$-to-PCA-64 embedding through systematic coding-scheme comparison and dimensionality analysis, and verifies its transfer to real EEG in Chapter~\ref{chap:arspi_net}.

Chapter~\ref{chap:arspi_net} is the empirical center of gravity of the dissertation. It validates the reservoir on real clinical EEG ($N = 211$), evaluates message-passing GNNs and alternative graph strategies, characterizes the graph propagation regime for small dense electrode graphs, and demonstrates that graph-topological analysis provides a clinically informative interpretability layer.

Chapter~\ref{chap:dynamics} develops the dynamical biomarker framework and tests hypotheses about how reservoir dynamics differ across affective and clinical states.

Chapter~\ref{chap:synthesis} synthesizes the temporal and graph-theoretic analyses into a unified structure--function coupling framework.

Chapter~\ref{chap:conclusion} consolidates contributions, outlines the boundaries of the present work, and proposes future directions including neuromorphic hardware deployment and single-trial EEG processing.

\subsection{Scope and Delimitations}

The dissertation contributes an electrical-engineering framework for extracting structured, interpretable representations from complex physiological time-series data. Its advances are concentrated in three areas: intrinsically interpretable neuromorphic signal processing, geometric characterization of nuisance structure in multichannel EEG, and clinically grounded structural analysis across multiple representational layers. The affective and clinical context is essential for validation, but the dissertation is not principally a psychology or neuroscience thesis. It is an engineering thesis about how to design, analyze, and validate a computational architecture that produces clinically meaningful, traceable characterizations of neural response organization from noisy neuroelectric measurements.

Several delimitations bound the scope. The interpretability claim is validated through a four-level experimental program (Chapters~\ref{chap:embeddings}--\ref{chap:synthesis}) and includes a direct comparison between ARSPI-Net's intrinsic temporal attention and EEGNet's post-hoc gradient saliency on the same data (Chapter~\ref{chap:arspi_net}). Formal XAI benchmarking with insertion/deletion curves, ROAR (remove-and-retrain), and clinician user studies following the ProtoPMed-EEG methodology \cite{barnett_protopmed_2024} are identified as future work. The energy-efficiency argument is motivational; the dissertation does not develop new neuromorphic hardware or measure energy consumption on physical chips. The spiking neuron model and learning framework are drawn from the existing literature rather than newly proposed. The clinical analysis adopts a transdiagnostic perspective across multiple diagnostic variables rather than targeting a single diagnostic boundary, and all clinical results are treated as hypothesis-generating pending independent replication. The reservoir-derived dynamical biomarkers are treated as calibrated functionals of a standardized nonlinear measurement instrument rather than as direct measures of cortical dynamics. The resulting framework is intended to support future low-power, interpretable biosignal systems for clinical decision support, brain--computer interfaces, and affective monitoring, while remaining grounded in the mathematics of dynamical systems, information extraction, and graph-based representation.